%-------------------------
% Cover Letter - Tavern Research Engineer I
% Author: Ajay Venkatesh
%-------------------------

\documentclass[letterpaper,11pt]{article}

\usepackage[top=0.8in, bottom=0.8in, left=0.9in, right=0.9in]{geometry}
\usepackage[hidelinks]{hyperref}
\usepackage{lmodern}
\usepackage{parskip}
\usepackage{microtype}

\setlength{\parskip}{6pt}
\setlength{\parindent}{0pt}

\begin{document}

\begin{flushleft}
\textbf{\large Ajay Venkatesh} \\
Brooklyn, NY \\
+1 516 585 9013 \\
\href{mailto:av3855@nyu.edu}{av3855@nyu.edu}
\end{flushleft}

\vspace{10pt}

May 13, 2026

\vspace{6pt}

Hiring Manager \\
Tavern Research

\vspace{10pt}

Dear Hiring Manager,

My name is Ajay Venkatesh --- finishing my M.S. in Computer Engineering at NYU, based here in Brooklyn, with several years of production software experience: a few years at PwC in Bengaluru, a summer SDE internship at Amazon, and ongoing on-campus work at NYU's Novel AI Technologies. My focus has been on backend systems, the practical AI layer --- LLM-powered features, RAG pipelines, and agentic workflows --- and shipping things end-to-end in small, fast-moving teams.

The Engineer I posting reads like a description of how I already work. \textbf{Python} is my daily driver for both backend and ML pipelines; \textbf{JavaScript / TypeScript} is where I spend my time on the product surface; I'm fluent in \textbf{Docker} and \textbf{Kubernetes} for containerization and deployment; and I've shipped APIs and database-backed systems across multiple stacks. At Amazon I deployed a \textbf{Retrieval-Augmented Generation (RAG)} pipeline with an \textbf{agentic workflow} that autonomously triages incidents by reasoning over operational logs, and implemented a \textbf{Model Context Protocol (MCP)} server for tool-augmented retrieval. At Campus Mesh I built \textbf{CampusIQ}, an LLM-powered AI assistant with iterative \textbf{prompt engineering} and \textbf{RAG retrieval} over structured data. Deploying LLMs and AI-assisted systems in production isn't a side project --- it's been the shape of my last year of work.

On using AI to accelerate development specifically: I use \textbf{Cursor} and \textbf{GitHub Copilot} daily, but I verify trade-offs and understand what's underneath before shipping production code. Fundamentals plus AI, not either alone --- which I think is what the Engineer I description is getting at when it pairs "ship clean code" with "use LLMs to expand your skill set."

On backend systems and reliability: at PwC I spent three years on a distributed backend on \textbf{Microsoft Azure} integrating with \textbf{Microsoft CRM}, with event-driven processing, fault-tolerant retries over SQL Server, and parallelized batch ingestion that cut execution from hours to minutes. At Amazon I provisioned \textbf{Infrastructure-as-Code} on \textbf{AWS CDK} (Lambda, API Gateway, DynamoDB, IAM, CloudWatch) with multi-stage CI/CD and rollback controls. I've debugged production incidents through logs, metrics, and traces --- the actual unglamorous work of keeping a system reliable.

What pulls me toward Tavern is the combination: a small team in Brooklyn, modern stack, AI-native workflows, and a product space where the work actually matters --- civic engagement isn't a synthetic problem. I'd like to work somewhere I can keep building at speed, contribute across the stack, and learn from senior engineers without the friction of a big-company process.

I'd welcome the chance to discuss further. Thanks for considering my application.

\vspace{12pt}

Sincerely, \\
Ajay Venkatesh

\end{document}
